\chapter{Clean up and next steps}
\label{ch:wrapup}

Congratulations! If you have followed this far, you have now implemented all components of LLM, from tokenizer to text generation. This chapter summarizes Volume 1 and introduces advanced topics to be covered in Volume 2,\textit{Make LLM-advanced}.

\section{Volume 1 Summary}

\subsection{implemented}

\begin{center}
\begin{tabularx}{\textwidth}{lX}
\toprule
\textbf{page}&\textbf{Implementation details}\\
\midrule
Chapter 3 & BPE, Unigram Tokenizer (Learning/Encoding/Decoding) \\
Chapter 4 & Token/Positional Embedding, EmbeddingStack \\
Chapter 5 & Scaled Dot-Product Attention, Causal Mask, Multi-Head Attention \\
Chapter 6 & FeedForward (GELU), Transformer Block (Pre-LN) \\
Chapter 7 & GPT Model, LMDataset, Cross-Entropy, AdamW, Cosine LR, Trainer \\
Chapter 8 & Greedy/Temperature/Top-K/Top-P Sampler, generate function \\
\bottomrule
\end{tabularx}
\end{center}

\subsection{test}

All 48 unit tests pass. We verified that all modules operate correctly through geometric tests, numerical tests, and integration tests.

\begin{lstlisting}[style=bashstyle]
$ pytest tests/ -v
========================= 48 passed in 2.47s =========================
\end{lstlisting}

\subsection{What the reader has}

Now the reader can:
\begin{itemize}
\item Learn tokenizer with your own corpus
\item Create a GPT model of arbitrary size (just change the config)
\item Run training loop on CPU/GPU
\item Generate text with various sampling strategies
\item Understand the inner workings of HuggingFace Transformers when reading their code
\end{itemize}

\section{Limitations of Volume 1}

To be honest, the model created in Volume 1 is far from ChatGPT. There are three main limitations:

\textbf{size}: Volume 1 model has 10 million$\sim$100 million parameters. GPT-3 has an estimated parameter of 175 billion, and GPT-4 has an estimated parameter of 1.7 trillion. Distributed learning is essential to overcome this size difference.

\textbf{architecture}: Volume 1 is close to the original Transformers from 2017. Modern LLM (LLaMA, Mistral, GPT-4) uses more efficient and stable components such as RoPE, GQA, SwiGLU, and RMSNorm.

\textbf{array}: Volume 1 model simply learned ``next token prediction''. To create a chatbot that answers user questions, alignment techniques such as SFT (Supervised Fine-Tuning), RLHF (Reinforcement Learning from Human Feedback), and DPO (Direct Preference Optimization) are required.

\section{Covered in Volume 2}

Volume 2\textit{Make LLM-advanced}deals with how to overcome the three limitations above.

\subsection{Distributed Learning (2  --  Chapter 4)}

\begin{itemize}
\item\textbf{DDP (Distributed Data Parallel)}: Replicate learning of the same model on multiple GPUs.
\item\textbf{FSDP (Fully Sharded Data Parallel)}: Save memory by splitting the model into shards.
\item\textbf{ZeRO-1/2/3}: Shading optimizer/gradient/parameters.
\item\textbf{Pipeline Parallel}: Place models on different GPUs for each layer.
\item\textbf{Tensor Parallel}: Splitting a single layer across multiple GPUs (Megatron-LM).
\end{itemize}

\subsection{Modern Architecture (5  --  Chapter 6)}

\begin{itemize}
\item\textbf{RoPE (Rotary Position Embedding)}: Position encoding with complex rotation. Excellent extrapolation.
\item\textbf{GQA (Grouped Query Attention)}: Save inference memory by reducing KV head.
\item\textbf{SwiGLU}: Improved FFN performance with GLU variant. Used by LLaMA.
\item\textbf{RMSNorm}: A lightweight version of LayerNorm.
\item\textbf{MoE (Mixture-of-Experts)}: Only the top K tokens for each token are activated in the expert network pool.
\end{itemize}

\subsection{Quantization and inference optimization (Chapter 7)}

\begin{itemize}
\item\textbf{INT8/INT4 Quantization}: Saves memory 4$\sim$8 times by converting weights to low-bit integers.
\item\textbf{AWQ}: Minimizing accuracy loss using activation value statistics.
\item\textbf{GPTQ}: Quantization using secondary information.
\item\textbf{KV Cache}: Accelerate inference by reusing K and V of previous tokens.
\item\textbf{PagedAttention}: Address memory fragmentation with page-level KV management in vLLM.
\end{itemize}

\subsection{Alignment and fine tuning (Chapter 8)}

\begin{itemize}
\item\textbf{SFT (Supervised Fine-Tuning)}: Learning chatbot behavior through supervised learning.
\item\textbf{RLHF}: Optimizing human preferences with compensation model + PPO.
\item\textbf{DPO}: Direct preference optimization without the complexity of RLHF.
\item\textbf{LoRA}: Efficient fine tuning with low-rank adapter.
\end{itemize}

\subsection{Data Pipeline (Chapter 9)}

\begin{itemize}
\item\textbf{quality filter}: Filter by length, language, and repetition rate.
\item\textbf{MinHash Deduplication}: Duplicate document detection using Jaccard similarity approximation.
\item\textbf{synthetic data}: Template-based learning data generation.
\end{itemize}

\section{Recommended Learning Path}

Recommended activities before moving on to Volume 2:

\begin{enumerate}
\item\textbf{small project}: Study texts in your field of interest (e.g. Wikipedia Korean, GitHub code) with Volume 1 code. About 100,000 tokens can be learned in 1 hour.
\item\textbf{read the paper}: Read ``Attention Is All You Need'' (Vaswani et al. 2017) and ``Language Models are Few-Shot Learners'' (Brown et al. 2020). Having read Volume 1, you will understand much better.
\item\textbf{HuggingFace read code}: Open\code{transformers/models/gpt2/modeling\_gpt2.py}. Compare how similar it is to the code we created and what differences there are.
\end{enumerate}

\section{finally}

LLM may seem complicated, but it is ultimately a combination of tokenizer, embedding, attention, FFN, and learning loop. If you learn these basics, you can understand any new architecture or new optimization technique. In Volume 2, we will build ChatGPT level technology on this foundation.

Never stop coding. Never stop experimenting. That's the only way to improve your skills.

\begin{flushright}
\textit{End of Volume 1}\\[0.5em]
\textbf{dirmich}
\end{flushright}
